\documentclass{article}

\usepackage{neurips_2026}

\usepackage[utf8]{inputenc}
\usepackage[T1]{fontenc}
\usepackage{hyperref}
\usepackage{url}
\usepackage{booktabs}
\usepackage{amsfonts}
\usepackage{amsmath}
\usepackage{nicefrac}
\usepackage{microtype}
\usepackage{xcolor}
\usepackage{graphicx}
\usepackage{algorithm}
\usepackage{algorithmic}

\title{AC-DPO: Adaptive Capacity Direct Preference Optimization}

\author{
  Aaron Chen \\
  Washington University in St. Louis \\
  \texttt{aaronc@wustl.edu} \\
  \And
  Sherlock Zhang \\
  Washington University in St. Louis \\
  \texttt{jiaxi.zhang@wustl.edu} \\
}

\begin{document}

\maketitle

\begin{abstract}
Direct Preference Optimization (DPO) typically employs a fixed-capacity adapter throughout training, regardless of the varying difficulty of preference pairs. We propose Adaptive Capacity DPO (AC-DPO), a two-stage training framework that aligns model capacity with data difficulty by combining curriculum learning with dynamic LoRA rank expansion. In Stage~1, the model trains on easy preference pairs using a low-rank adapter ($r{=}8$, 295K parameters), then merges the learned weights into the base model. In Stage~2, a high-rank adapter ($r{=}64$, 2.36M parameters) is initialized for hard preference pairs. We evaluate AC-DPO against three baselines on a 10,000-pair subset of the Anthropic RLHF dataset using GPT-2. While all methods achieve comparable overall accuracy (43--44\%), AC-DPO attains the highest token-level reward margin on hard tasks (+18\% over the fixed $r{=}64$ baseline), confirming that capacity--difficulty alignment improves fine-grained discrimination on complex reasoning pairs. Ablation studies with a reverse curriculum and a fixed $r{=}8$ baseline further reveal that curriculum ordering interacts with capacity allocation in non-trivial ways, providing insights for future work on larger models.
\end{abstract}

%------------------------------------------------------------------
\section{Introduction}
%------------------------------------------------------------------

Aligning large language models (LLMs) with human preferences is a central challenge in modern NLP. Direct Preference Optimization \citep{rafailov2023dpo} offers an efficient alternative to reinforcement learning from human feedback (RLHF) by directly optimizing a policy from preference pairs without training a separate reward model. In practice, DPO is often combined with parameter-efficient fine-tuning methods such as Low-Rank Adaptation (LoRA) \citep{hu2022lora}, which inserts trainable low-rank matrices into frozen pretrained weights.

A key limitation of standard DPO with LoRA is that the adapter rank---and thus the model's trainable capacity---remains fixed throughout training. However, preference datasets are inherently heterogeneous: some pairs exhibit large quality gaps between chosen and rejected responses (``easy'' pairs), while others require subtle reasoning to distinguish (``hard'' pairs). Training a high-capacity adapter on easy data risks overfitting to superficial patterns, while a low-capacity adapter may lack the expressiveness needed for hard reasoning tasks.

Curriculum learning \citep{bengio2009curriculum} addresses data heterogeneity by ordering training examples from easy to hard. Recent work has applied curriculum strategies to DPO \citep{pattnaik2024currydpo}, but these methods adjust only the \emph{data} schedule while keeping model capacity fixed. Separately, dynamic rank methods have been explored in other domains \citep{deng2026drlora, croitoru2026curriculumdpo}, but not in the context of LLM preference optimization.

We propose \textbf{Adaptive Capacity DPO (AC-DPO)}, which jointly adapts both the data curriculum and the model's parameter capacity. Our contributions are:
\begin{enumerate}
    \item A two-stage training pipeline that combines curriculum-based data ordering with dynamic LoRA rank expansion, using adapter merging for knowledge transfer between stages.
    \item An automated difficulty partitioning scheme using a pretrained reward model (DeBERTa-v3) to score and split preference pairs by reward margin.
    \item A comprehensive evaluation including two ablation studies (reverse curriculum and fixed low-rank baseline) that isolate the effects of curriculum ordering and capacity expansion.
\end{enumerate}

%------------------------------------------------------------------
\section{Related Works}
%------------------------------------------------------------------

\paragraph{Direct Preference Optimization.}
DPO \citep{rafailov2023dpo} reformulates the RLHF objective by deriving a closed-form mapping between the reward function and the optimal policy, eliminating the need for a separate reward model during training. Given a preference dataset of triples $(\mathbf{x}, \mathbf{y}_w, \mathbf{y}_l)$ where $\mathbf{y}_w$ is preferred over $\mathbf{y}_l$, DPO directly optimizes the policy $\pi_\theta$ against a reference policy $\pi_{\text{ref}}$.

\paragraph{Low-Rank Adaptation.}
LoRA \citep{hu2022lora} parameterizes weight updates as $\Delta W = BA$ where $B \in \mathbb{R}^{d \times r}$ and $A \in \mathbb{R}^{r \times k}$ with rank $r \ll \min(d, k)$. This reduces trainable parameters by orders of magnitude while preserving model quality. The choice of rank $r$ controls the capacity--efficiency trade-off: higher ranks enable more expressive updates but increase computational cost and overfitting risk.

\paragraph{Curriculum Learning.}
Curriculum learning \citep{bengio2009curriculum} trains models on progressively harder examples. Curry-DPO \citep{pattnaik2024currydpo} applies this principle to preference optimization by ranking pairs by difficulty and training in easy-to-hard order, demonstrating improved alignment quality. However, Curry-DPO uses a fixed model capacity throughout training.

\paragraph{Dynamic Model Capacity.}
DR-LoRA \citep{deng2026drlora} introduces dynamic rank expansion for mixture-of-experts in text-to-image generation, progressively increasing LoRA rank during training. Curriculum-DPO++ \citep{croitoru2026curriculumdpo} combines data and model curricula for diffusion models. Our work extends this concept to LLM preference optimization, where the interaction between adapter capacity and preference pair difficulty has not been studied.

%------------------------------------------------------------------
\section{Background \& Preliminaries}
%------------------------------------------------------------------

\subsection{Direct Preference Optimization}

Given a dataset $\mathcal{D} = \{(\mathbf{x}^{(i)}, \mathbf{y}_w^{(i)}, \mathbf{y}_l^{(i)})\}_{i=1}^{N}$ of prompts with preferred and dispreferred responses, DPO minimizes:
\begin{equation}
\mathcal{L}_{\text{DPO}}(\pi_\theta; \pi_{\text{ref}}) = -\mathbb{E}_{(\mathbf{x}, \mathbf{y}_w, \mathbf{y}_l) \sim \mathcal{D}} \left[ \log \sigma \left( \beta \log \frac{\pi_\theta(\mathbf{y}_w | \mathbf{x})}{\pi_{\text{ref}}(\mathbf{y}_w | \mathbf{x})} - \beta \log \frac{\pi_\theta(\mathbf{y}_l | \mathbf{x})}{\pi_{\text{ref}}(\mathbf{y}_l | \mathbf{x})} \right) \right]
\label{eq:dpo}
\end{equation}
where $\sigma$ is the sigmoid function and $\beta$ controls the deviation from the reference policy $\pi_{\text{ref}}$.

\subsection{Low-Rank Adaptation}

For a pretrained weight matrix $W_0 \in \mathbb{R}^{d \times k}$, LoRA introduces a low-rank update:
\begin{equation}
W = W_0 + \Delta W = W_0 + BA
\end{equation}
where $B \in \mathbb{R}^{d \times r}$, $A \in \mathbb{R}^{r \times k}$, and $r$ is the adapter rank. The number of trainable parameters per adapted layer is $r(d + k)$, compared to $dk$ for full fine-tuning.

\subsection{Difficulty Scoring}

We define the difficulty of a preference pair $(\mathbf{x}, \mathbf{y}_w, \mathbf{y}_l)$ using the reward margin from a pretrained reward model $R$:
\begin{equation}
m(\mathbf{x}, \mathbf{y}_w, \mathbf{y}_l) = R(\mathbf{x}, \mathbf{y}_w) - R(\mathbf{x}, \mathbf{y}_l)
\label{eq:margin}
\end{equation}
Large positive margins indicate easy pairs where the quality difference is obvious; small or negative margins indicate hard pairs requiring nuanced judgment.

%------------------------------------------------------------------
\section{Method}
%------------------------------------------------------------------

AC-DPO consists of three phases: difficulty-based data partitioning, low-capacity training on easy data, and high-capacity training on hard data with knowledge transfer.

\subsection{Data Partitioning}

Given a preference dataset $\mathcal{D}$, we compute the reward margin $m_i$ for each pair using a pretrained reward model (Equation~\ref{eq:margin}). We sort pairs by margin in descending order and split at the median into:
\begin{align}
\mathcal{D}_{\text{easy}} &= \{(\mathbf{x}, \mathbf{y}_w, \mathbf{y}_l) \in \mathcal{D} \mid m \geq m_{\text{median}}\} \\
\mathcal{D}_{\text{hard}} &= \{(\mathbf{x}, \mathbf{y}_w, \mathbf{y}_l) \in \mathcal{D} \mid m < m_{\text{median}}\}
\end{align}

\subsection{Two-Stage Training with Capacity Expansion}

\paragraph{Stage 1: Easy curriculum with low rank.}
We attach a LoRA adapter with rank $r_1$ (small) to the base model $\theta_0$ and train on $\mathcal{D}_{\text{easy}}$ using the DPO objective (Equation~\ref{eq:dpo}). The low rank constrains the model to learn broad, high-signal patterns without overfitting.

\paragraph{Knowledge Transfer via Adapter Merging.}
After Stage~1, we permanently merge the learned low-rank weights into the base model:
\begin{equation}
\theta_1 = \theta_0 + B_1 A_1
\end{equation}
This preserves the knowledge from Stage~1 while freeing the adapter slot for Stage~2.

\paragraph{Stage 2: Hard curriculum with high rank.}
We initialize a new LoRA adapter with rank $r_2 \gg r_1$ on the merged model $\theta_1$ and train on $\mathcal{D}_{\text{hard}}$. The higher rank provides the capacity needed for complex reasoning patterns, while the merged base already encodes basic alignment from Stage~1.

The complete procedure is summarized in Algorithm~\ref{alg:acdpo}.

\begin{algorithm}[t]
\caption{AC-DPO: Adaptive Capacity Direct Preference Optimization}
\label{alg:acdpo}
\begin{algorithmic}[1]
\REQUIRE Base model $\theta_0$, reference model $\pi_{\text{ref}}$, preference data $\mathcal{D}$, reward model $R$, ranks $r_1 < r_2$
\STATE Compute margins $m_i = R(\mathbf{x}_i, \mathbf{y}_{w,i}) - R(\mathbf{x}_i, \mathbf{y}_{l,i})$ for all $i$
\STATE Split $\mathcal{D}$ into $\mathcal{D}_{\text{easy}}$ (high margin) and $\mathcal{D}_{\text{hard}}$ (low margin) at median
\STATE Initialize LoRA adapter $(B_1, A_1)$ with rank $r_1$ on $\theta_0$
\STATE Train on $\mathcal{D}_{\text{easy}}$ with DPO loss $\rightarrow$ obtain trained $(B_1, A_1)$
\STATE Merge: $\theta_1 \leftarrow \theta_0 + B_1 A_1$
\STATE Initialize LoRA adapter $(B_2, A_2)$ with rank $r_2$ on $\theta_1$
\STATE Train on $\mathcal{D}_{\text{hard}}$ with DPO loss $\rightarrow$ obtain trained $(B_2, A_2)$
\STATE Merge: $\theta_{\text{final}} \leftarrow \theta_1 + B_2 A_2$
\RETURN $\theta_{\text{final}}$
\end{algorithmic}
\end{algorithm}



%------------------------------------------------------------------
\section{Experiments}
%------------------------------------------------------------------

\subsection{Experimental Setup}

\paragraph{Dataset.}
We use a 10,000-pair subset from the Anthropic RLHF dataset \citep{bai2022training} (\texttt{Dahoas/rm-static}). Each pair consists of a prompt with a preferred (chosen) and dispreferred (rejected) response. We score all pairs using the DeBERTa-v3-large reward model \citep{he2021deberta} (\texttt{OpenAssistant/reward-model-deberta-v3-large-v2}) and split at the median reward margin into 5,000 easy pairs (margin $\geq 0.572$) and 5,000 hard pairs (margin $< 0.572$). The easy subset has a mean margin of 1.698, while the hard subset has a mean margin of $-0.151$, indicating that many hard pairs have negative margins where the reward model disagrees with the human label. Each subset is further split 80/20 into train (4,000) and evaluation (1,000) sets with a fixed seed of 5100.

\paragraph{Base Model.}
We use GPT-2 (124M parameters) \citep{radford2019gpt2} as the base model, with a frozen copy serving as the reference model $\pi_{\text{ref}}$. All LoRA adapters target the \texttt{c\_attn} (combined query-key-value) projection layers.

\paragraph{Training Configuration.}
All experiments use $\beta = 0.1$, batch size 4, 4 training epochs, and the AdamW optimizer. AC-DPO Stage~1 uses learning rate $10^{-4}$ (rank $r{=}8$, 294,912 trainable parameters, 0.24\% of total); Stage~2 uses learning rate $5 \times 10^{-5}$ (rank $r{=}64$, 2,359,296 trainable parameters, 1.86\% of total). Sequences are truncated to 512 tokens with a maximum prompt length of 256 tokens. All experiments use the same random seed for reproducibility.

\paragraph{Compared Methods.}
We evaluate four configurations to isolate the effects of curriculum ordering and capacity expansion:
\begin{itemize}
    \item \textbf{AC-DPO} (ours): Easy$\rightarrow$Hard curriculum with $r{=}8 \rightarrow r{=}64$ expansion.
    \item \textbf{Baseline $r{=}64$}: Fixed $r{=}64$ on shuffled mixed data (standard DPO + LoRA).
    \item \textbf{Baseline $r{=}8$}: Fixed $r{=}8$ on shuffled mixed data (low-capacity control).
    \item \textbf{Reverse Curriculum}: Hard$\rightarrow$Easy with $r{=}8 \rightarrow r{=}64$ (ablation for curriculum direction).
\end{itemize}

\paragraph{Evaluation Metrics.}
For each evaluation pair, we compute the log-probability of the chosen and rejected responses under the trained model. We report:
\begin{itemize}
    \item \textbf{Sum Accuracy}: fraction of pairs where $\sum \log p(\mathbf{y}_w | \mathbf{x}) > \sum \log p(\mathbf{y}_l | \mathbf{x})$.
    \item \textbf{Token Accuracy}: fraction where the average per-token log-probability favors the chosen response.
    \item \textbf{Token Margin}: mean difference in per-token log-probabilities, $\bar{\log p}(\mathbf{y}_w) - \bar{\log p}(\mathbf{y}_l)$, measuring the model's confidence in its preference.
\end{itemize}

\subsection{Main Results}

Table~\ref{tab:main_results} presents the evaluation results across all four methods.

\begin{table}[t]
\caption{Evaluation results on 1,000 easy, 1,000 hard, and 2,000 combined preference pairs. Best results per column are \textbf{bolded}. AC-DPO achieves the highest token margin on hard tasks, while all methods perform comparably on overall accuracy.}
\label{tab:main_results}
\centering
\small
\begin{tabular}{l cc cc cc}
\toprule
& \multicolumn{2}{c}{\textbf{Easy}} & \multicolumn{2}{c}{\textbf{Hard}} & \multicolumn{2}{c}{\textbf{All}} \\
\cmidrule(lr){2-3} \cmidrule(lr){4-5} \cmidrule(lr){6-7}
\textbf{Method} & Sum Acc & Tok Acc & Sum Acc & Tok Acc & Sum Acc & Tok Acc \\
\midrule
AC-DPO (ours)      & 34.7 & 56.4 & \textbf{51.8} & \textbf{52.6} & 43.25 & 54.5 \\
Baseline $r{=}64$  & 35.4 & 58.9 & 51.6 & 52.1 & 43.5  & 55.5 \\
Baseline $r{=}8$   & 35.7 & 61.0 & 51.4 & 52.1 & 43.55 & 56.55 \\
Reverse Curriculum  & \textbf{36.1} & \textbf{62.8} & 51.7 & 52.7 & \textbf{43.9} & \textbf{57.75} \\
\bottomrule
\end{tabular}
\end{table}

\begin{table}[t]
\caption{Token-level reward margins (higher is better) and training efficiency. AC-DPO achieves the strongest discrimination on hard pairs while using fewer trainable parameters in Stage~1.}
\label{tab:margins}
\centering
\small
\begin{tabular}{l ccc ccc}
\toprule
& \multicolumn{3}{c}{\textbf{Token Margin}} & \multicolumn{3}{c}{\textbf{Training Efficiency}} \\
\cmidrule(lr){2-4} \cmidrule(lr){5-7}
\textbf{Method} & Easy & Hard & All & Params & Runtime (s) & GPU (GB) \\
\midrule
AC-DPO (ours)      & 0.146 & \textbf{0.059} & 0.103 & 295K$\rightarrow$2.36M & 2047 & 9.06 \\
Baseline $r{=}64$  & 0.236 & 0.050 & 0.143 & 2.36M & 2061 & 9.06 \\
Baseline $r{=}8$   & 0.285 & 0.047 & 0.166 & 295K & 2131 & 9.02 \\
Reverse Curriculum  & \textbf{0.331} & 0.038 & \textbf{0.185} & 295K$\rightarrow$2.36M & 2122 & 9.06 \\
\bottomrule
\end{tabular}
\end{table}

\paragraph{Overall accuracy is comparable across methods.}
All four methods achieve overall sum accuracy in the narrow range of 43.25--43.9\%, suggesting that at the GPT-2 scale with 10K training pairs, the choice of curriculum ordering and adapter rank has limited impact on aggregate preference accuracy. This is consistent with the observation that GPT-2's base capacity may be the primary bottleneck.

\paragraph{AC-DPO achieves the strongest hard-task discrimination.}
On hard preference pairs, AC-DPO attains the highest token margin (0.059), exceeding the fixed $r{=}64$ baseline by 18\% (0.050) and the fixed $r{=}8$ baseline by 26\% (0.047). This indicates that dedicating high capacity specifically to hard data---after first learning basic patterns with low capacity---produces sharper per-token distinctions on complex reasoning pairs.

\paragraph{Easy-task performance trades off with hard-task focus.}
AC-DPO shows lower easy-task token accuracy (56.4\%) compared to baselines (58.9--62.8\%). This is expected: Stage~1 uses only $r{=}8$ (295K parameters) on easy data, providing less capacity than the $r{=}64$ baselines. The merged Stage~1 knowledge may also be partially overwritten during Stage~2 training on hard data.

\subsection{Ablation Studies}

\paragraph{Effect of curriculum direction (Reverse Curriculum).}
The reverse curriculum (Hard$\rightarrow$Easy) achieves the highest overall accuracy (43.9\%) and token accuracy (57.75\%), outperforming AC-DPO's easy-to-hard ordering. This initially surprising result can be explained by \emph{recency bias}: the model's final performance is disproportionately influenced by the last training stage. The reverse curriculum dedicates its high-capacity Stage~2 to easy data, which has clearer preference signals and constitutes half the evaluation set. However, its hard-task token margin (0.038) is the lowest among all methods, confirming that it sacrifices hard-task discrimination for easy-task performance---the opposite trade-off from AC-DPO.

\paragraph{Effect of adapter capacity (Baseline $r{=}8$).}
The fixed $r{=}8$ baseline performs comparably to the $r{=}64$ baseline on sum accuracy (43.55\% vs 43.5\%) and even outperforms it on token accuracy (56.55\% vs 55.5\%). This suggests that at the GPT-2 scale, the low-rank adapter may provide beneficial regularization. However, the $r{=}8$ baseline has the lowest hard-task token margin (0.047), indicating that while low capacity suffices for coarse preference accuracy, it limits fine-grained discrimination on difficult pairs.

\subsection{Training Dynamics}

Figure~\ref{fig:training} illustrates the training loss trajectories. AC-DPO Stage~1 converges rapidly on easy data (final loss 0.464), reflecting the clear preference signals. Stage~2 starts with elevated loss ($\sim$1.1) due to the newly initialized high-rank adapter and harder data, then decreases to 0.715. The baseline trains on mixed data with a smoother trajectory (final loss 0.624). The reverse curriculum shows the opposite pattern: Stage~1 on hard data converges slowly (loss 0.572), while Stage~2 on easy data converges faster (loss 0.638).

\begin{figure}[t]
\centering
\fbox{\parbox{0.9\linewidth}{\centering\vspace{2em}
\textit{Training loss curves: AC-DPO (Stage~1 + Stage~2), Baseline $r{=}64$, Reverse Curriculum (Stage~1 + Stage~2). AC-DPO Stage~1 converges fastest on easy data; Stage~2 starts high due to new adapter initialization on hard data.}
\vspace{2em}}}
\caption{Training loss over steps. AC-DPO Stage~1 (easy, $r{=}8$) converges rapidly to 0.464. Stage~2 (hard, $r{=}64$) begins at $\sim$1.1 after adapter re-initialization and decreases to 0.715. The baseline on mixed data converges to 0.624.}
\label{fig:training}
\end{figure}

\subsection{Compute Efficiency Analysis}

Table~\ref{tab:compute} compares the effective parameter-compute across methods. We define parameter-steps as the product of trainable parameters and training steps, measuring the total optimization work.

\begin{table}[h]
\caption{Compute efficiency comparison. AC-DPO uses 55\% of the parameter-steps of the fixed $r{=}64$ baseline by leveraging low-rank training in Stage~1.}
\label{tab:compute}
\centering
\small
\begin{tabular}{l rrr}
\toprule
\textbf{Method} & \textbf{Param-Steps ($\times 10^9$)} & \textbf{Relative} & \textbf{Wall Time (s)} \\
\midrule
AC-DPO (ours)     & $0.295\text{M} \times 4000 + 2.36\text{M} \times 4000 = 10.6$ & 0.56$\times$ & 2047 \\
Baseline $r{=}64$ & $2.36\text{M} \times 8000 = 18.9$ & 1.00$\times$ & 2061 \\
Baseline $r{=}8$  & $0.295\text{M} \times 8000 = 2.4$ & 0.13$\times$ & 2131 \\
Reverse Curriculum & $0.295\text{M} \times 4000 + 2.36\text{M} \times 4000 = 10.6$ & 0.56$\times$ & 2122 \\
\bottomrule
\end{tabular}
\end{table}

AC-DPO uses only 56\% of the parameter-steps of the $r{=}64$ baseline while achieving comparable accuracy and superior hard-task discrimination. The wall-clock times are similar because GPU utilization is dominated by forward/backward passes through the full model rather than the adapter parameters alone. On larger models where adapter computation becomes a larger fraction of total cost, the parameter-step savings would translate more directly to wall-clock savings.



%------------------------------------------------------------------
\section{Conclusions}
%------------------------------------------------------------------

We presented AC-DPO, a framework that aligns model capacity with data difficulty during preference optimization by combining curriculum learning with dynamic LoRA rank expansion. Our two-stage pipeline---low-rank training on easy pairs followed by high-rank training on hard pairs with adapter merging for knowledge transfer---achieves the highest token-level reward margin on hard preference pairs (+18\% over the fixed-rank baseline) while using only 56\% of the parameter-compute budget. Although overall accuracy differences are small at the GPT-2 scale, the consistent advantage in hard-task discrimination validates the core hypothesis that matching capacity to difficulty improves fine-grained preference learning.

Our ablation studies reveal two additional insights: (1)~curriculum direction matters---the reverse curriculum trades hard-task discrimination for easy-task performance, confirming that the easy-to-hard ordering in AC-DPO specifically benefits complex reasoning; (2)~at the GPT-2 scale, low-rank adapters provide implicit regularization that can match high-rank adapters on coarse accuracy, but not on fine-grained token-level margins.

%------------------------------------------------------------------
\section{Limitations \& Future Work}
%------------------------------------------------------------------

\paragraph{Model scale.}
Our experiments use GPT-2 (124M parameters), where the base model's capacity may be the primary bottleneck, compressing all methods into a narrow accuracy range. The benefits of dynamic capacity allocation are expected to be more pronounced on larger models (e.g., LLaMA-7B or 13B) where the gap between $r{=}8$ and $r{=}64$ represents a more meaningful fraction of the model's expressiveness.

\paragraph{Binary difficulty split.}
We partition data into only two difficulty levels (easy/hard) at the median. A finer-grained, multi-stage curriculum with gradual rank expansion (e.g., $r{=}8 \rightarrow 16 \rightarrow 32 \rightarrow 64$) could provide smoother capacity--difficulty alignment and reduce the loss spike observed at stage transitions.

\paragraph{Catastrophic forgetting.}
AC-DPO's lower easy-task performance suggests that Stage~2 training partially overwrites Stage~1 knowledge despite adapter merging. Techniques such as experience replay (mixing easy examples into Stage~2) or elastic weight consolidation could mitigate this forgetting.

\paragraph{Evaluation scope.}
We evaluate only log-probability-based preference accuracy. Future work should include generation-quality metrics (e.g., win rates from GPT-4 judges, human evaluation) and downstream task performance to assess whether improved token margins translate to better generated outputs.

\paragraph{Recency bias.}
The reverse curriculum's strong overall performance highlights a recency bias in two-stage training: the final stage disproportionately influences evaluation metrics. Investigating training schedules that interleave difficulty levels or use replay buffers could address this limitation.

%------------------------------------------------------------------
\begin{ack}
We thank the instructors and TAs of CSE 5100 (Spring 2026) at Washington University in St. Louis for their guidance. Experiments were conducted on the WashU Academic Compute Cluster using NVIDIA A40 GPUs allocated through the \texttt{condo-cse5100} partition.
\end{ack}

%------------------------------------------------------------------
\bibliographystyle{plainnat}

\begin{thebibliography}{10}

\bibitem[Bai et~al.(2022)]{bai2022training}
Bai, Y., Jones, A., Ndousse, K., Askell, A., Chen, A., DasSarma, N., Drain, D., Fort, S., Ganguli, D., Henighan, T., et~al.
\newblock Training a helpful and harmless assistant with reinforcement learning from human feedback.
\newblock \emph{arXiv preprint arXiv:2204.05862}, 2022.

\bibitem[Bengio et~al.(2009)]{bengio2009curriculum}
Bengio, Y., Louradour, J., Collobert, R., and Weston, J.
\newblock Curriculum learning.
\newblock In \emph{Proceedings of the 26th International Conference on Machine Learning (ICML)}, pp.~41--48, 2009.

\bibitem[Croitoru et~al.(2024)]{croitoru2024curriculum}
Croitoru, F.-A., Hondru, V., Ionescu, R.~T., Sebe, N., and Shah, M.
\newblock Curriculum direct preference optimization for diffusion and consistency models.
\newblock \emph{arXiv preprint arXiv:2405.13637}, 2024.

\bibitem[Croitoru et~al.(2026)]{croitoru2026curriculumdpo}
Croitoru, F.-A., Hondru, V., Ionescu, R.~T., Sebe, N., and Shah, M.
\newblock Curriculum-DPO++: Direct preference optimization via data and model curricula for text-to-image generation.
\newblock \emph{arXiv preprint arXiv:2602.13055}, 2026.

\bibitem[Deng et~al.(2026)]{deng2026drlora}
Deng, G., Li, B., Chen, R., Wang, H., Wen, L., and Song, L.
\newblock DR-LoRA: Dynamic rank LoRA for mixture-of-experts adaptation.
\newblock \emph{arXiv preprint arXiv:2601.04823}, 2026.

\bibitem[He et~al.(2021)]{he2021deberta}
He, P., Gao, J., and Chen, W.
\newblock DeBERTaV3: Improving DeBERTa using ELECTRA-style pre-training with gradient-disentangled embedding sharing.
\newblock \emph{arXiv preprint arXiv:2111.09543}, 2021.

\bibitem[Hu et~al.(2022)]{hu2022lora}
Hu, E.~J., Shen, Y., Wallis, P., Allen-Zhu, Z., Li, Y., Wang, S., Wang, L., and Chen, W.
\newblock LoRA: Low-rank adaptation of large language models.
\newblock In \emph{International Conference on Learning Representations (ICLR)}, 2022.

\bibitem[Pattnaik et~al.(2024)]{pattnaik2024currydpo}
Pattnaik, P., Maheshwary, R., Ogueji, K., Yadav, V., and Madhusudhan, S.~T.
\newblock Curry-DPO: Enhancing alignment using curriculum learning \& ranked preferences.
\newblock \emph{arXiv preprint arXiv:2403.07230}, 2024.

\bibitem[Radford et~al.(2019)]{radford2019gpt2}
Radford, A., Wu, J., Child, R., Luan, D., Amodei, D., and Sutskever, I.
\newblock Language models are unsupervised multitask learners.
\newblock \emph{OpenAI Technical Report}, 2019.

\bibitem[Rafailov et~al.(2023)]{rafailov2023dpo}
Rafailov, R., Sharma, A., Mitchell, E., Ermon, S., Manning, C.~D., and Finn, C.
\newblock Direct preference optimization: Your language model is secretly a reward model.
\newblock In \emph{Advances in Neural Information Processing Systems (NeurIPS)}, 2023.

\end{thebibliography}

%%%%%%%%%%%%%%%%%%%%%%%%%%%%%%%%%%%%%%%%%%%%%%%%%%%%%%%%%%%%
\appendix

\section{Detailed Per-Step Training Metrics}

Table~\ref{tab:training_dynamics} reports training metrics at key checkpoints for all methods.

\begin{table}[h]
\caption{Training loss and reward accuracy at selected steps.}
\label{tab:training_dynamics}
\centering
\small
\begin{tabular}{ll rrr}
\toprule
\textbf{Method} & \textbf{Checkpoint} & \textbf{Loss} & \textbf{Reward Acc} & \textbf{Margin} \\
\midrule
AC-DPO Stage~1 & Step 1000 & 0.580 & 0.775 & 0.733 \\
($r{=}8$, easy) & Step 2000 & 0.420 & 0.750 & 1.620 \\
                & Step 4000 & 0.597 & 0.725 & 1.246 \\
\midrule
AC-DPO Stage~2 & Step 1000 & 0.748 & 0.525 & $-$0.038 \\
($r{=}64$, hard) & Step 2000 & 0.686 & 0.475 & 0.053 \\
                 & Step 4000 & 0.663 & 0.650 & 0.112 \\
\midrule
Baseline $r{=}64$ & Step 2000 & 0.617 & 0.625 & 0.268 \\
(mixed)            & Step 4000 & 0.684 & 0.600 & 0.116 \\
                   & Step 8000 & 0.680 & 0.475 & 0.171 \\
\midrule
Reverse Stage~1 & Step 2000 & 0.669 & 0.575 & 0.070 \\
($r{=}8$, hard)  & Step 4000 & 0.572 & 0.775 & 0.341 \\
\midrule
Reverse Stage~2 & Step 2000 & 0.475 & 0.750 & 0.959 \\
($r{=}64$, easy) & Step 4000 & 0.639 & 0.625 & 0.661 \\
\midrule
Baseline $r{=}8$ & Step 2000 & 0.618 & 0.650 & 0.322 \\
(mixed)           & Step 4000 & 0.738 & 0.600 & 0.116 \\
                  & Step 8000 & 0.769 & 0.500 & 0.145 \\
\bottomrule
\end{tabular}
\end{table}

\section{Experimental Infrastructure}

All experiments were conducted on a single NVIDIA A40 GPU (48GB VRAM) on the WashU Academic Compute Cluster. Peak GPU memory usage was approximately 9.06~GB across all configurations. The complete experimental pipeline (data preparation, four training runs, and evaluation) required approximately 2.5 hours of GPU time. Software stack: Python 3.10, PyTorch, Hugging Face Transformers, PEFT, and TRL.

%%%%%%%%%%%%%%%%%%%%%%%%%%%%%%%%%%%%%%%%%%%%%%%%%%%%%%%%%%%%

\newpage
\input{checklist.tex}

\end{document}
